%% Research design section 12 -- the experiment design, and section 12's own
%% success criteria.
%%
%% Every table in this section is \input from tables/, generated by
%% `make tables` from experiments/results/*.md. A table typed by hand here is a
%% number whose provenance banner was dropped on the way in.
%%
%% Each table's note carries the banner of the md it came from:
%%   MOCK          - not performance numbers
%%   REAL SIM      - LLMServingSim, named cache dir, not hardware
%%   REAL HARDWARE - node serials named
%%
%% Experiments, in the order they are reported:
%%   E-G1  / E-G1b  compression and top-K against a binding spec (mock)
%%   E-G3           the same correctness check under the real simulator
%%   E-G4           the contention model against a microbenchmark (hardware)
%%   E-G5           deployment on real hardware, recommendation vs boundary
%%   E-G6           scalability, 32-128 devices by symmetry
%%   E-G7           holdout, ablation, baseline fairness
%%
%% The two correctness numbers are reported for every arm, in every table, and
%% never aggregated away.
\section{Evaluation}
\label{sec:eval}

\subsection{Setup}

Three evaluators are used and never confused. A deterministic \emph{mock} obeys
the same physics as the elimination bounds, which is what makes a disagreement
between the search and an exhaustive oracle meaningful; it is not a
performance model and no figure produced by it is a measurement.
\emph{LLMServingSim} with an analytical network backend is the real simulator.
\emph{Hardware} is an eight-GPU node whose accelerator serial numbers every raw
file records. Every table carries its provenance in the caption.

Every fixture in the accompanying repository is fictional and says so in its
header, with one exception: the hardware cluster description is generated from
the microbenchmark's own raw files, so that each figure marked as measured is
traceable to a file recording the node's serials, the process binding, the load
average and which other processes held a GPU.

The success criteria, the numeric limits, and the interval a threshold was
chosen from were registered before the corresponding experiment ran, and the
registration is append-only. Where a limit was derived from a measurement, the
runs it was derived from are named and excluded from the set that tests it.

\subsection{Correctness under the real simulator}

% claim: C10
% claim: C11
% claim: C12
\input{e_g3_real_sim_oracle}

Table~\ref{tab:e_g3_real_sim_oracle} reports the search against an exhaustive
oracle, both driven by LLMServingSim. Both correctness numbers are zero on all
three fixtures: no feasible candidate was eliminated, and no two placements the
oracle prices differently were folded together.

% claim: C13
The \texttt{complete} column is \textbf{False} on all three, and that
qualification travels with the result everywhere it is quoted. Between six and
eight per cent of placements failed to simulate at all. They are reported as
\texttt{unknown\_measurement}---the evaluator was asked and did not
answer---and not as infeasible. The correctness claim is therefore about the
placements the simulator judged, which is a narrower statement than it would be
if those failures had been counted as rejections. Their cause is not
established and is listed as an open item in Section~\ref{sec:limits}.
\pending{P1.5: SIM\_ERROR cause}

\input{e_g3_real_sim_oracle_2}

Table~\ref{tab:e_g3_real_sim_oracle_2} reports what the compression cost
against what it saved. On the largest fixture, \egthreeSimOracleAbcde~seconds
of oracle simulation become \egthreeSimProposedAbcde, a saving of
\egthreeSavingAbcde~seconds with the hashing, the isomorphism checks and the
bound evaluation all charged against the compression rather than ignored.

The compression ratio is a property of the graph and not of the evaluator: it
agrees to four decimals between the mock and the real simulator on every shared
fixture, which is what a structural claim should do when the thing that
evaluates it is replaced.

\subsection{Compression, and where it fails}

% claim: C1
% claim: C4
\input{e_g1_toy_pilot}
\input{e_g6_scale}

Table~\ref{tab:e_g1_toy_pilot} reports the compression on the toy corpus.
Table~\ref{tab:e_g6_scale} reports it against device count and cluster
symmetry, for a mock evaluator, at 32, 64 and 128 devices.

% claim: C23
The symmetry axis carries the result. At 128 devices the same
\egsixEmbeddingsLarge{} embeddings compress to \egsixRepsAsym{} representatives
on a fully asymmetric cluster and to \egsixRepsSym{} on a fully symmetric one.
The asymmetric ratio is \egsixRatioAsym---that is, the compression finds almost
nothing to fold, and the corpus reports that rather than omitting the
configuration. A method whose benefit depends on a property of the input should
be shown on inputs that lack it.

This experiment has no success target and is registered as report-only, with a
single failure condition: the compression must not cost more than it saves on a
fully symmetric cluster. It does not. The grid is driven by the mock, and a
simulation at this scale is not offered as validation of accuracy at scale.

\subsection{Contribution A: the boundary in the signature}

% claim: C24
\input{e_g7_ablation}

Table~\ref{tab:e_g7_ablation} reports the ablation. With the shared boundary in
the equivalence signature the corpus yields no mis-merged pairs. With
\texttt{no\_boundary}---every other label retained---two placements that differ
only in which uplink they cross are folded into one representative, and the
oracle prices them apart.

The threshold that decides ``prices them apart'' is derived from the oracle's
own distribution rather than fixed in advance. An earlier form of this
experiment used a fixed tight specification and found no mis-merges, because
that specification made both placements infeasible and the oracle agreed they
were equally bad. A mis-merge count depends on where the threshold sits, so the
derivation is stated with the count.

\subsection{Search quality}

% claim: C5
% claim: C6
\input{e_g1b_topk}
\input{e_g2_topk_holdout}
\begin{figure}[t]\centering
  \includegraphics[width=\linewidth]{topk_holdout.pdf}
  \caption{Feasible recall against the simulation budget, both arms, on the holdout fixtures~\sourcetag{MOCK}. The oracle's level is the dotted line. At $K=4$ the graph search is behind, which the curve shows without the reader comparing two rows.}
  \label{fig:topk}
\end{figure}

Table~\ref{tab:e_g1b_topk} and Table~\ref{tab:e_g2_topk_holdout} report
feasible recall against the existing planner's surrogate top-$K$ under a
specification whose SLOs bind. At $K = 8$ and $K = 16$ this search reaches
higher recall. At $K = 4$ it does not---it ties on one holdout fixture and
trails on the other---and that is the row to read before the others.

The baseline is scored generously and deliberately so. It ranks and judges
templates rather than placements, so it cannot name which devices a candidate
occupies; one verdict is therefore credited to every placement of the template
it judged. There is no stricter reading that would be fair to it, and winning
on a technicality would not be winning.

% claim: C25
% claim: C26
\input{e_g7_holdout}
\input{e_g7_baseline_fairness}

Table~\ref{tab:e_g7_holdout} reports a holdout fixed before the scaling grid
ran: \egsevenHoldoutEmbeddings{} embeddings compressing to
\egsevenHoldoutReps{} representatives, both correctness numbers zero, and
higher recall than the surrogate at the registered operating point. The second
holdout---a hardware fixture with one uplink reservation altered---does not
exist, because the hardware campaign has not produced it. That row is reported
as not run rather than omitted, since an omitted row reads as a row that
passed. \pending{E-G5: real-lab holdout}

\subsection{The contention model against the wire}

% claim: C15
% claim: C16
\input{e_g4_microbench_2}
\begin{figure}[t]\centering
  \includegraphics[width=\linewidth]{contention_error.pdf}
  \caption{Fluid and null error against message size, on the declaration the measurements support~\sourcetag{REAL HARDWARE}. The bidirectional error is near zero from 4 to 64 MiB and rises at 128: an accuracy domain, not a single figure.}
  \label{fig:contention}
\end{figure}

Table~\ref{tab:e_g4_microbench_2} reports the registered verdict, and it is a
\textbf{failure}: on the topology the plan declared before measuring,
processor-sharing is \egfourFluidPlannedSame\,\% out at the median on the
condition it exists for, against \egfourNullPlannedSame\,\% for pricing each
flow alone.

The reason is not the model. Two device pairs believed to share a host bridge
were measured to share nothing: each concurrent copy sustained the full
single-copy rate, and a third process saturating one pair left the other at its
unloaded rate to three digits. Run on the declaration the measurements support,
the same model is \egfourFluidMeasuredSame\,\% out; and on the bidirectional
condition, where the two flows genuinely do share endpoints, it is
\egfourFluidMeasuredBidi\,\% against \egfourNullMeasuredBidi\,\% for the null
model.

Both are reported because they are different corrections. One would rewrite the
contention model; the other rewrites a line of a cluster description. The
registered verdict is computed on the registered declaration, and the corrected
figure is not offered in its place.

% claim: C22
% claim: C29
\paragraph{The same question at a second location, with the opposite answer.}
Across an inter-node NIC between two machines, two concurrent streams take
exactly half each: a single stream reaches \egfourNicSingle~Gbit/s and two
reach \egfourNicTwoSame~each. That is processor sharing precisely, and it is
what the intra-node PCIe pairs did not do.

The pair of results is worth more than either alone. It is not that the model
is right in one place and wrong in another---it is that \emph{which resource
is genuinely shared} is a claim about hardware, settled by measurement at each
location separately, and the model prices whatever it is told to. The
bidirectional case says the same thing twice: the intra-node pair reaches only
a third again as much as a single direction, while the NIC exceeds twice it,
because that wire really is full duplex.

Two of the five registered conditions remain unanswered at this location, for
reasons that do not merge. Two concurrent flows over \emph{disjoint} NICs is
impossible on this hardware---each machine has exactly one InfiniBand
device---while the collective is deferred, needing an NCCL environment the peer
node does not yet have. \pending{E-G4: collective at location (b)}

% claim: C17
The model's accuracy domain has an upper end. Above a message size the hardware
begins overlapping the two directions of one link, the pair exceeds what
processor sharing predicts, and the model becomes optimistic by about a third.
That boundary is stated rather than absorbed into a fitted capacity: a capacity
chosen because it reproduces the answer is not a measurement, and it would make
the next node's prediction wrong silently.

\subsection{Hardware}

% pending
This section reports the deployment campaign: the recommended placement against
its SLOs, the boundary alternatives against the bounds that ranked or rejected
them, and the difference between two placements of one template across a shared
interconnect. The campaign is registered in the accompanying repository and has
not completed. \pending{E-G5: the hardware matrix}

% claim: C21
% pending
One result from it is already measured and is reported here because it does not
depend on the rest. Two placements of a single template---a tensor-parallel
pair on an NVLink link, and the same pair across a PCIe bridge---differ in
served $p99$ time to first token by a factor that grows sharply at the knee,
because the slower interconnect moves the knee itself. Three repetitions per
cell; the ranges do not overlap. The planner being extended cannot express the
difference between these two placements at all, which is the point.
